%%%%%%%%%%%%%%%%%%%%%%%%%%%%%%%%%%%%%%%%%%%%%%%%%%%%%%%%%%%%%%%%%%%%%%%%%%%%%%%
%% app_gauge_nlpp.tex
%%
%% Appendix: Gauge Transformation of the Nonlocal Pseudopotential and
%%           Explicit Evaluation of [r_alpha, V^nl] for KB Pseudopotentials
%%
%% %%%%%%%%%%%%%%%%%%%%%%%%%%%%%%%%%%%%%%%%%%%%%%%%%%%%%%%%%%%%%%%%%%%%%%%%%%%%%%%
%% app_gauge_nlpp.tex
%%
%% Appendix: Gauge Transformation of the Nonlocal Pseudopotential and
%%           Explicit Evaluation of [r_alpha, V^nl] for KB Pseudopotentials
%%
%% %%%%%%%%%%%%%%%%%%%%%%%%%%%%%%%%%%%%%%%%%%%%%%%%%%%%%%%%%%%%%%%%%%%%%%%%%%%%%%%
%% app_gauge_nlpp.tex
%%
%% Appendix: Gauge Transformation of the Nonlocal Pseudopotential and
%%           Explicit Evaluation of [r_alpha, V^nl] for KB Pseudopotentials
%%
%% %%%%%%%%%%%%%%%%%%%%%%%%%%%%%%%%%%%%%%%%%%%%%%%%%%%%%%%%%%%%%%%%%%%%%%%%%%%%%%%
%% app_gauge_nlpp.tex
%%
%% Appendix: Gauge Transformation of the Nonlocal Pseudopotential and
%%           Explicit Evaluation of [r_alpha, V^nl] for KB Pseudopotentials
%%
%% \input{app_gauge_nlpp} goes in the appendix section of paper2_main.tex.
%%
%% Structure:
%%   A. Gauge transformation V^nl under LG -> VG: derives the coupling
%%      -i A_alpha [r_alpha, V^nl] that appears in H^VG
%%   B. Explicit algebra of [r_alpha, V^nl] for single-projector KB PP
%%   C. Generalization to multiple projectors per channel
%%   D. Nonlocal current density J_nl from the commutator
%%   E. Connection to AppNls_0xyz() in RMG
%%
%% CONVENTIONS: same as sec_theory_current.tex (atomic units, e_electron=-1)
%% KEY SOURCE: NEAT-project-notes_main.tex lines 785-851 (Starace derivation)
%%             Kleinman-Bylander (1982); Mattiat-Luber (2022) eq.12
%%%%%%%%%%%%%%%%%%%%%%%%%%%%%%%%%%%%%%%%%%%%%%%%%%%%%%%%%%%%%%%%%%%%%%%%%%%%%%%

\section{Gauge Transformation of the Nonlocal Pseudopotential
  and the Commutator $[\hat{r}_\alpha, \Vnl]$}
\label{app:gauge_nlpp}

This appendix derives two connected results:
(A)~the transformation of the nonlocal Kleinman-Bylander (KB)
pseudopotential from the length gauge to the velocity gauge,
establishing why the commutator $[\hat{r}_\alpha, \Vnl]$ enters the
velocity gauge Hamiltonian; and (B)~the explicit algebraic evaluation
of the commutator for a separable pseudopotential, providing the
formula implemented in \texttt{AppNls\_0xyz()} in \texttt{RMG}.

%--------------------------------------------------------------------
\subsection*{A.\; Gauge Transformation of $\Vnl$ from Length to
  Velocity Gauge}
\label{app:gauge-transform-Vnl}
%--------------------------------------------------------------------

Under the gauge transformation from length gauge (LG) to velocity
gauge (VG), the wavefunction and all operators transform by the
unitary operator\cite{Luber-rt-tddft,Mattiat-Luber-2022}
\begin{equation}
  U(\mathbf{r},t)
  = \exp\!\left(\imath\frac{q}{\hbar c}\chi(\mathbf{r},t)\right)
  \stackrel{\mathrm{a.u.}}{=}
  \exp(-\imath\chi(\mathbf{r},t)),
  \label{eq:U-gauge}
\end{equation}
where in atomic units with $\hbar = c = 1$ and $q = -1$ for
electrons, $q/\hbar c = -1$.
The gauge function for the LG $\to$ VG transformation is
\begin{equation}
  \chi(\mathbf{r},t) = -\mathbf{r}\cdot\Avec(t),
  \label{eq:chi-vg}
\end{equation}
where $\Avec(t) = -\int_0^t \mathbf{E}(t')\,\mathrm{d}t'$ is the
vector potential chosen so that $\phi'=0$ in the velocity
gauge.\cite{Luber-rt-tddft,rt-DFTB-Wong-2023}
With $\chi = -\mathbf{r}\cdot\Avec$, the transformation operator is
\begin{equation}
  U(\mathbf{r},t)
  = e^{-\imath(-\mathbf{r}\cdot\Avec)}
  = e^{\imath\mathbf{r}\cdot\Avec(t)}.
  \label{eq:U-explicit}
\end{equation}

Any operator $\hat{O}$ transforms as $\hat{O}' = U\hat{O}U^\dag$.
For a \emph{local} operator $f(\mathbf{r})$, the transformation is
trivial: $U f(\mathbf{r})U^\dag = f(\mathbf{r})$, because
$U(\mathbf{r})$ commutes with any function of $\mathbf{r}$.

For the \emph{nonlocal} pseudopotential $\Vnl$, the two operators
$U$ and $\Vnl$ do not commute, and the transformed potential is
\begin{equation}
  \tilde{V}^{\mathrm{nl}}
  = U\Vnl U^\dag
  = e^{\imath\mathbf{r}\cdot\Avec}\Vnl
    e^{-\imath\mathbf{r}\cdot\Avec}.
  \label{eq:Vnl-transformed}
\end{equation}
To extract the first-order coupling to $\Avec$, we use the
Baker-Campbell-Hausdorff (BCH) expansion truncated at first order
in $\Avec$ (valid within linear response):
\begin{align}
  e^{\imath\mathbf{r}\cdot\Avec}\Vnl e^{-\imath\mathbf{r}\cdot\Avec}
  &= \Vnl
   + \imath[\mathbf{r}\cdot\Avec,\, \Vnl]
   + \mathcal{O}(\Avec^2)
  \nonumber \\
  &= \Vnl
   + \imath A_\alpha(t)[\hat{r}_\alpha,\Vnl]
   + \mathcal{O}(\Avec^2),
  \label{eq:Vnl-BCH}
\end{align}
where repeated Cartesian index $\alpha$ is summed.
The velocity gauge Hamiltonian therefore differs from the length
gauge Hamiltonian by the first-order correction
\begin{equation}
  H^{\mathrm{VG}} - H^{\mathrm{LG}}
  \supset \imath A_\alpha(t)[\hat{r}_\alpha,\Vnl]
  + \mathcal{O}(\Avec^2).
  \label{eq:H-VG-correction}
\end{equation}
This correction is the origin of the nonlocal pseudopotential term
in the velocity gauge Hamiltonian.
When combined with the paramagnetic coupling
$-\Avec\cdot\hat{\mathbf{p}}$ (from the kinetic energy), the
total field-electron coupling in the velocity gauge is
\begin{equation}
  V_{\mathrm{field}}^{\mathrm{VG}}
  = -A_\alpha(t)\hat{p}_\alpha
    + \imath A_\alpha(t)[\hat{r}_\alpha,\Vnl]
  = A_\alpha(t)\!\left(-\hat{p}_\alpha
    + \imath[\hat{r}_\alpha,\Vnl]\right),
  \label{eq:V-field-VG}
\end{equation}
which can be written as $-A_\alpha(t)\hat{v}_\alpha$ upon
identifying $\hat{v}_\alpha = \hat{p}_\alpha -
\imath[\hat{r}_\alpha,\Vnl] = \hat{p}_\alpha +
\frac{1}{\imath}[\hat{r}_\alpha,\Vnl]\cdot(-1)$.
(The sign is consistent with $\hat{v}_\alpha = \hat{p}_\alpha +
\frac{1}{\imath}[\hat{r}_\alpha,\Vnl]$ at $A=0$, as derived in
Section~\ref{sec:velocity-eom}.)

%--------------------------------------------------------------------
\subsection*{B.\; Explicit Evaluation of $[\hat{r}_\alpha, \Vnl]$
  for Kleinman-Bylander Pseudopotentials}
\label{app:KB-commutator}
%--------------------------------------------------------------------

The Kleinman-Bylander separable pseudopotential for atom $I$ in the
single-projector-per-channel approximation
is\cite{kleinman-bylander-1982}
\begin{equation}
  \Vnl = \sum_I\sum_{l,m}
  |\beta^{lm}_I\rangle\, h^l_I\,
  \langle\beta^{lm}_I|,
  \label{eq:KB-def}
\end{equation}
where $|\beta^{lm}_I\rangle$ is the KB projector for atom $I$
with angular momentum quantum numbers $(l,m)$, centered at
$\mathbf{R}_I$, and $h^l_I$ is the (real, scalar) strength
coefficient for angular momentum channel $l$.
In position representation, the projector wavefunction is
$\beta^{lm}_I(\mathbf{r}) = \langle\mathbf{r}|\beta^{lm}_I\rangle$.
The sum over $(l,m)$ runs over all angular momentum channels and
magnetic quantum numbers included in the pseudopotential.

We evaluate the commutator $[\hat{r}_\alpha, \Vnl]$ by direct
computation, acting on an arbitrary state $|\psi\rangle$:
\begin{align}
  [\hat{r}_\alpha, \Vnl]|\psi\rangle
  &= \hat{r}_\alpha\left(\Vnl|\psi\rangle\right)
   - \Vnl\left(\hat{r}_\alpha|\psi\rangle\right).
  \label{eq:comm-expand}
\end{align}

\paragraph{First term: $\hat{r}_\alpha(\Vnl|\psi\rangle)$.}
Inserting the KB form of $\Vnl$:
\begin{align}
  \hat{r}_\alpha\left(\Vnl|\psi\rangle\right)
  &= \hat{r}_\alpha\sum_{I,l,m}
     |\beta^{lm}_I\rangle\, h^l_I\,
     \langle\beta^{lm}_I|\psi\rangle
  \nonumber \\
  &= \sum_{I,l,m}
     \hat{r}_\alpha|\beta^{lm}_I\rangle\, h^l_I\,
     \langle\beta^{lm}_I|\psi\rangle,
  \label{eq:term1}
\end{align}
where $\hat{r}_\alpha$ acts on the ket $|\beta^{lm}_I\rangle$
(multiplying the position wavefunction $\beta^{lm}_I(\mathbf{r})$
by $r_\alpha$), while the scalar overlap
$c_{I,lm} \equiv \langle\beta^{lm}_I|\psi\rangle
= \int\beta^{lm*}_I(\mathbf{r}')\psi(\mathbf{r}')\,\mathrm{d}^3r'$
is unaffected.

\paragraph{Second term: $\Vnl(\hat{r}_\alpha|\psi\rangle)$.}
\begin{align}
  \Vnl\left(\hat{r}_\alpha|\psi\rangle\right)
  &= \sum_{I,l,m}
     |\beta^{lm}_I\rangle\, h^l_I\,
     \langle\beta^{lm}_I|\hat{r}_\alpha|\psi\rangle.
  \label{eq:term2}
\end{align}
Since $\hat{r}_\alpha$ is Hermitian,
$\langle\beta^{lm}_I|\hat{r}_\alpha|\psi\rangle
= \langle\hat{r}_\alpha\beta^{lm}_I|\psi\rangle
= \int r_\alpha\,\beta^{lm*}_I(\mathbf{r})\,\psi(\mathbf{r})\,
\mathrm{d}^3r$.
We denote this as $\langle r_\alpha\beta^{lm}_I|\psi\rangle$,
where $r_\alpha\beta^{lm}_I$ means the function
$r_\alpha\beta^{lm*}_I(\mathbf{r})$ in position space.

\paragraph{The commutator.}
Subtracting eq.~(\ref{eq:term2}) from eq.~(\ref{eq:term1}):
\begin{align}
  [\hat{r}_\alpha,\Vnl]|\psi\rangle
  &= \sum_{I,l,m} h^l_I
     \Big[
     \hat{r}_\alpha|\beta^{lm}_I\rangle\,\langle\beta^{lm}_I|\psi\rangle
   - |\beta^{lm}_I\rangle\,\langle\hat{r}_\alpha\beta^{lm}_I|\psi\rangle
     \Big].
  \label{eq:comm-intermediate}
\end{align}
Using the shorthand $|r_\alpha\beta^{lm}_I\rangle \equiv
\hat{r}_\alpha|\beta^{lm}_I\rangle$ for the position-weighted
projector ket, this is
\begin{equation}
  \boxed{
  [\hat{r}_\alpha,\Vnl]|\psi\rangle
  = \sum_{I,l,m} h^l_I\!\left[
    |r_\alpha\beta^{lm}_I\rangle\langle\beta^{lm}_I|\psi\rangle
  - |\beta^{lm}_I\rangle\langle r_\alpha\beta^{lm}_I|\psi\rangle
    \right].
  }
  \label{eq:comm-result}
\end{equation}
This is a rank-2 update (in the language of matrix algebra): the
commutator of a separable operator with $\hat{r}_\alpha$ is
again separable, but now involves two types of projectors — one
weighted by $r_\alpha$ in the ket, one in the bra.

Equation~(\ref{eq:comm-result}) has a transparent structure:
\begin{itemize}
  \item The \emph{first term}
    $|r_\alpha\beta^{lm}_I\rangle\langle\beta^{lm}_I|\psi\rangle$
    keeps the unmodified projector in the bra and weights the ket
    by $r_\alpha$.
    Physically: the electron is projected onto the original KB
    projector at $\mathbf{r}'$ and then the result is multiplied
    by $r_\alpha$ at the field point $\mathbf{r}$.
  \item The \emph{second term}
    $|\beta^{lm}_I\rangle\langle r_\alpha\beta^{lm}_I|\psi\rangle$
    keeps the unmodified projector ket but weights the bra by
    $r_\alpha$.
    Physically: the electron is projected onto the
    $r_\alpha$-weighted projector (measuring overlap with the
    position-weighted projector) and the result propagates with
    the original projector shape.
\end{itemize}
The two terms have opposite signs and together ensure that the
commutator operator is anti-Hermitian, consistent with
$[\hat{r}_\alpha, \Vnl]^\dag = -[\hat{r}_\alpha, \Vnl]$ for
Hermitian $\hat{r}_\alpha$ and $\Vnl$.

\paragraph{Matrix element between Kohn-Sham orbitals.}
Taking the matrix element of eq.~(\ref{eq:comm-result}) between
$\langle\psi_a|$ and $|\psi_b\rangle$:
\begin{align}
  \langle\psi_a|[\hat{r}_\alpha,\Vnl]|\psi_b\rangle
  &= \sum_{I,l,m} h^l_I
     \Big[
       \underbrace{\langle\psi_a|r_\alpha\beta^{lm}_I\rangle}_{
         \displaystyle A^{lm}_{I,a}}
       \underbrace{\langle\beta^{lm}_I|\psi_b\rangle}_{
         \displaystyle B^{lm}_{I,b}}
     - \underbrace{\langle\psi_a|\beta^{lm}_I\rangle}_{
         \displaystyle \overline{B^{lm}_{I,a}}}
       \underbrace{\langle r_\alpha\beta^{lm}_I|\psi_b\rangle}_{
         \displaystyle \overline{A^{lm}_{I,b}}}
     \Big],
  \label{eq:comm-matrix-element}
\end{align}
where we introduced the notation
\begin{align}
  B^{lm}_{I,b}
  &= \langle\beta^{lm}_I|\psi_b\rangle
   = \int \beta^{lm*}_I(\mathbf{r})\,\psi_b(\mathbf{r})
     \,\mathrm{d}^3r,
  \label{eq:B-overlap}
  \\
  A^{lm}_{I,b}
  &= \langle r_\alpha\beta^{lm}_I|\psi_b\rangle
   = \int r_\alpha\,\beta^{lm*}_I(\mathbf{r})\,\psi_b(\mathbf{r})
     \,\mathrm{d}^3r.
  \label{eq:A-overlap}
\end{align}
Here $A^{lm}_{I,b}$ is the overlap of the $\alpha$-position-weighted
projector $r_\alpha\beta^{lm}_I$ with the orbital $\psi_b$, and
$B^{lm}_{I,b}$ is the standard KB projector overlap.
Both are computed as spatial integrals over the real-space grid.

Equation~(\ref{eq:comm-matrix-element}) therefore requires two types
of integrals per projector per orbital: the standard KB overlap
$B^{lm}_{I,b}$ (already available from the usual nonlocal PP force
evaluation) and the new position-weighted overlap $A^{lm}_{I,b}$.

%--------------------------------------------------------------------
\subsection*{C.\; Generalization to Multiple Projectors per Channel}
\label{app:multi-projector}
%--------------------------------------------------------------------

In the general KB form with multiple projectors per angular momentum
channel (eq.~11 of Kleinman and Bylander\cite{kleinman-bylander-1982}),
\begin{equation}
  \Vnl = \sum_I\sum_{l,m}\sum_{i,j}
  |\beta^{lm}_{I,i}\rangle\, h^l_{I,ij}\,
  \langle\beta^{lm}_{I,j}|,
  \label{eq:KB-general}
\end{equation}
where $i,j$ are radial projector indices and $h^l_{I,ij}$ is the
(generally non-diagonal) strength matrix.
The commutator result of eq.~(\ref{eq:comm-result}) extends directly
to
\begin{equation}
  [\hat{r}_\alpha,\Vnl]|\psi\rangle
  = \sum_{I,l,m}\sum_{i,j} h^l_{I,ij}\!\left[
    |r_\alpha\beta^{lm}_{I,i}\rangle
    \langle\beta^{lm}_{I,j}|\psi\rangle
  - |\beta^{lm}_{I,i}\rangle
    \langle r_\alpha\beta^{lm}_{I,j}|\psi\rangle
    \right].
  \label{eq:comm-result-general}
\end{equation}
The algebraic structure is identical; only the matrix of strength
coefficients $h^l_{I,ij}$ becomes non-diagonal in the projector
indices.
For the norm-conserving pseudopotentials used in \texttt{RMG}, the
single-projector case (eq.~\ref{eq:comm-result}) applies.

%--------------------------------------------------------------------
\subsection*{D.\; Nonlocal Contribution to the Current Density}
\label{app:Jnl-expression}
%--------------------------------------------------------------------

The nonlocal current density contribution to
eq.~(\ref{eq:J-total}) is defined through
eq.~(\ref{eq:Jnl-continuity}).
We now verify that this definition is consistent with the velocity
operator result by showing that the volume integral of
$\mathbf{J}_\mathrm{nl}$ gives the expected macroscopic contribution.

Using the single-particle density matrix $P(\mathbf{r},\mathbf{r}')
= \sum_n f_n \psi_n(\mathbf{r})\psi_n^*(\mathbf{r}')$, the
nonlocal contribution to the macroscopic current along direction
$\alpha$ is
\begin{align}
  J_{\mathrm{nl},\alpha}(t)
  &= \int_\Omega\!
     \operatorname{Re}\!\left[
       \psi^*(\mathbf{r},t)\,
       \frac{1}{\imath}[\hat{r}_\alpha,\Vnl]\,
       \psi(\mathbf{r},t)
     \right]\mathrm{d}^3r
  \nonumber \\
  &= \frac{1}{\imath}\sum_{I,l,m} h^l_I\,
     \operatorname{Re}\!\Bigg[
       \int\!\psi^*(\mathbf{r})\,r_\alpha\,\beta^{lm}_I(\mathbf{r})
       \,\mathrm{d}^3r
       \cdot
       \int\!\beta^{lm*}_I(\mathbf{r}')\,\psi(\mathbf{r}')
       \,\mathrm{d}^3r'
  \nonumber \\
  &\qquad\qquad\qquad\quad
   - \int\!\psi^*(\mathbf{r})\,\beta^{lm}_I(\mathbf{r})
     \,\mathrm{d}^3r
     \cdot
     \int\!r_\alpha\,\beta^{lm*}_I(\mathbf{r}')\,\psi(\mathbf{r}')
     \,\mathrm{d}^3r'
     \Bigg].
  \label{eq:Jnl-full}
\end{align}
In the density matrix notation of Section~\ref{sec:current-dm-sec},
this is evaluated as the matrix element
$\langle\psi_a|\frac{1}{\imath}[\hat{r}_\alpha,\Vnl]|\psi_b\rangle$
from eq.~(\ref{eq:comm-matrix-element}), contracted with $P_{ab}(t)$.
No additional real-space integration is required at each time step
beyond what is already computed in building $\Jmat{\alpha}_{ab}$.

%--------------------------------------------------------------------
\subsection*{E.\; Connection to \texttt{AppNls\_0xyz()} in RMG}
\label{app:code-connection}
%--------------------------------------------------------------------

In \texttt{RMG}, the commutator $[\hat{r}_\alpha,\Vnl]|\psi_b\rangle$
is evaluated by the function \texttt{AppNls\_0xyz()}, called from
\texttt{CurrentNlpp()} for directions $\alpha \in \{x,y,z\}$.

The function computes exactly eq.~(\ref{eq:comm-result}) on the
real-space grid.
For each orbital $\psi_b$ and each atom $I$, it evaluates:
\begin{enumerate}
  \item The standard projector overlap
    $B^{lm}_{I,b} = \langle\beta^{lm}_I|\psi_b\rangle$
    (a sum over grid points in the support of $\beta^{lm}_I$);
  \item The position-weighted overlap
    $A^{lm}_{I,b} = \langle r_\alpha\beta^{lm}_I|\psi_b\rangle$
    (the same sum but with each grid point contribution weighted
    by its $r_\alpha$ coordinate relative to atom $I$);
  \item The output array as
    $[\hat{r}_\alpha,\Vnl]\psi_b(\mathbf{r})
    = \sum_{I,l,m} h^l_I\!\left[
      r_\alpha\beta^{lm}_I(\mathbf{r})\cdot B^{lm}_{I,b}
      - \beta^{lm}_I(\mathbf{r})\cdot A^{lm}_{I,b}
    \right]$,
    stored in the work array \texttt{nv[]}.
\end{enumerate}
The result is passed to \texttt{CurrentNlpp()}, which computes the
matrix element $\langle\psi_a|\texttt{nv}\rangle$ via a BLAS
\texttt{gemm} operation, multiplies by \texttt{I\_t} $= \imath$,
and adds the result to \texttt{Palpha\_matrix}.
As shown in Section~\ref{sec:imath-convention}, the combined factor
$\imath \times \frac{1}{\imath} = 1$ ensures that the nonlocal PP
correction enters the total current matrix
$\boldsymbol{\mathcal{J}}^\alpha$ with unit weight, consistent with
eq.~(\ref{eq:Jmat-full}).

A self-consistency check: at the ground state, when all orbitals
$\psi_n$ are real-valued eigenstates, the density matrix is
$P_{ab} = f_a\delta_{ab}$ and the imaginary part of
$\langle\psi_a|\frac{1}{\imath}[\hat{r}_\alpha,\Vnl]|\psi_a\rangle$
should vanish because $[\hat{r}_\alpha,\Vnl]$ is anti-Hermitian
(the diagonal elements of an anti-Hermitian operator are purely
imaginary for complex orbitals and zero for real orbitals).
Thus $J_\mathrm{nl}(t=0) = 0$, consistent with vanishing ground-state
current.
This provides a simple numerical test that can be run on the
ground-state density matrix before starting the RT-TDDFT propagation.
 goes in the appendix section of paper2_main.tex.
%%
%% Structure:
%%   A. Gauge transformation V^nl under LG -> VG: derives the coupling
%%      -i A_alpha [r_alpha, V^nl] that appears in H^VG
%%   B. Explicit algebra of [r_alpha, V^nl] for single-projector KB PP
%%   C. Generalization to multiple projectors per channel
%%   D. Nonlocal current density J_nl from the commutator
%%   E. Connection to AppNls_0xyz() in RMG
%%
%% CONVENTIONS: same as sec_theory_current.tex (atomic units, e_electron=-1)
%% KEY SOURCE: NEAT-project-notes_main.tex lines 785-851 (Starace derivation)
%%             Kleinman-Bylander (1982); Mattiat-Luber (2022) eq.12
%%%%%%%%%%%%%%%%%%%%%%%%%%%%%%%%%%%%%%%%%%%%%%%%%%%%%%%%%%%%%%%%%%%%%%%%%%%%%%%

\section{Gauge Transformation of the Nonlocal Pseudopotential
  and the Commutator $[\hat{r}_\alpha, \Vnl]$}
\label{app:gauge_nlpp}

This appendix derives two connected results:
(A)~the transformation of the nonlocal Kleinman-Bylander (KB)
pseudopotential from the length gauge to the velocity gauge,
establishing why the commutator $[\hat{r}_\alpha, \Vnl]$ enters the
velocity gauge Hamiltonian; and (B)~the explicit algebraic evaluation
of the commutator for a separable pseudopotential, providing the
formula implemented in \texttt{AppNls\_0xyz()} in \texttt{RMG}.

%--------------------------------------------------------------------
\subsection*{A.\; Gauge Transformation of $\Vnl$ from Length to
  Velocity Gauge}
\label{app:gauge-transform-Vnl}
%--------------------------------------------------------------------

Under the gauge transformation from length gauge (LG) to velocity
gauge (VG), the wavefunction and all operators transform by the
unitary operator\cite{Luber-rt-tddft,Mattiat-Luber-2022}
\begin{equation}
  U(\mathbf{r},t)
  = \exp\!\left(\imath\frac{q}{\hbar c}\chi(\mathbf{r},t)\right)
  \stackrel{\mathrm{a.u.}}{=}
  \exp(-\imath\chi(\mathbf{r},t)),
  \label{eq:U-gauge}
\end{equation}
where in atomic units with $\hbar = c = 1$ and $q = -1$ for
electrons, $q/\hbar c = -1$.
The gauge function for the LG $\to$ VG transformation is
\begin{equation}
  \chi(\mathbf{r},t) = -\mathbf{r}\cdot\Avec(t),
  \label{eq:chi-vg}
\end{equation}
where $\Avec(t) = -\int_0^t \mathbf{E}(t')\,\mathrm{d}t'$ is the
vector potential chosen so that $\phi'=0$ in the velocity
gauge.\cite{Luber-rt-tddft,rt-DFTB-Wong-2023}
With $\chi = -\mathbf{r}\cdot\Avec$, the transformation operator is
\begin{equation}
  U(\mathbf{r},t)
  = e^{-\imath(-\mathbf{r}\cdot\Avec)}
  = e^{\imath\mathbf{r}\cdot\Avec(t)}.
  \label{eq:U-explicit}
\end{equation}

Any operator $\hat{O}$ transforms as $\hat{O}' = U\hat{O}U^\dag$.
For a \emph{local} operator $f(\mathbf{r})$, the transformation is
trivial: $U f(\mathbf{r})U^\dag = f(\mathbf{r})$, because
$U(\mathbf{r})$ commutes with any function of $\mathbf{r}$.

For the \emph{nonlocal} pseudopotential $\Vnl$, the two operators
$U$ and $\Vnl$ do not commute, and the transformed potential is
\begin{equation}
  \tilde{V}^{\mathrm{nl}}
  = U\Vnl U^\dag
  = e^{\imath\mathbf{r}\cdot\Avec}\Vnl
    e^{-\imath\mathbf{r}\cdot\Avec}.
  \label{eq:Vnl-transformed}
\end{equation}
To extract the first-order coupling to $\Avec$, we use the
Baker-Campbell-Hausdorff (BCH) expansion truncated at first order
in $\Avec$ (valid within linear response):
\begin{align}
  e^{\imath\mathbf{r}\cdot\Avec}\Vnl e^{-\imath\mathbf{r}\cdot\Avec}
  &= \Vnl
   + \imath[\mathbf{r}\cdot\Avec,\, \Vnl]
   + \mathcal{O}(\Avec^2)
  \nonumber \\
  &= \Vnl
   + \imath A_\alpha(t)[\hat{r}_\alpha,\Vnl]
   + \mathcal{O}(\Avec^2),
  \label{eq:Vnl-BCH}
\end{align}
where repeated Cartesian index $\alpha$ is summed.
The velocity gauge Hamiltonian therefore differs from the length
gauge Hamiltonian by the first-order correction
\begin{equation}
  H^{\mathrm{VG}} - H^{\mathrm{LG}}
  \supset \imath A_\alpha(t)[\hat{r}_\alpha,\Vnl]
  + \mathcal{O}(\Avec^2).
  \label{eq:H-VG-correction}
\end{equation}
This correction is the origin of the nonlocal pseudopotential term
in the velocity gauge Hamiltonian.
When combined with the paramagnetic coupling
$-\Avec\cdot\hat{\mathbf{p}}$ (from the kinetic energy), the
total field-electron coupling in the velocity gauge is
\begin{equation}
  V_{\mathrm{field}}^{\mathrm{VG}}
  = -A_\alpha(t)\hat{p}_\alpha
    + \imath A_\alpha(t)[\hat{r}_\alpha,\Vnl]
  = A_\alpha(t)\!\left(-\hat{p}_\alpha
    + \imath[\hat{r}_\alpha,\Vnl]\right),
  \label{eq:V-field-VG}
\end{equation}
which can be written as $-A_\alpha(t)\hat{v}_\alpha$ upon
identifying $\hat{v}_\alpha = \hat{p}_\alpha -
\imath[\hat{r}_\alpha,\Vnl] = \hat{p}_\alpha +
\frac{1}{\imath}[\hat{r}_\alpha,\Vnl]\cdot(-1)$.
(The sign is consistent with $\hat{v}_\alpha = \hat{p}_\alpha +
\frac{1}{\imath}[\hat{r}_\alpha,\Vnl]$ at $A=0$, as derived in
Section~\ref{sec:velocity-eom}.)

%--------------------------------------------------------------------
\subsection*{B.\; Explicit Evaluation of $[\hat{r}_\alpha, \Vnl]$
  for Kleinman-Bylander Pseudopotentials}
\label{app:KB-commutator}
%--------------------------------------------------------------------

The Kleinman-Bylander separable pseudopotential for atom $I$ in the
single-projector-per-channel approximation
is\cite{kleinman-bylander-1982}
\begin{equation}
  \Vnl = \sum_I\sum_{l,m}
  |\beta^{lm}_I\rangle\, h^l_I\,
  \langle\beta^{lm}_I|,
  \label{eq:KB-def}
\end{equation}
where $|\beta^{lm}_I\rangle$ is the KB projector for atom $I$
with angular momentum quantum numbers $(l,m)$, centered at
$\mathbf{R}_I$, and $h^l_I$ is the (real, scalar) strength
coefficient for angular momentum channel $l$.
In position representation, the projector wavefunction is
$\beta^{lm}_I(\mathbf{r}) = \langle\mathbf{r}|\beta^{lm}_I\rangle$.
The sum over $(l,m)$ runs over all angular momentum channels and
magnetic quantum numbers included in the pseudopotential.

We evaluate the commutator $[\hat{r}_\alpha, \Vnl]$ by direct
computation, acting on an arbitrary state $|\psi\rangle$:
\begin{align}
  [\hat{r}_\alpha, \Vnl]|\psi\rangle
  &= \hat{r}_\alpha\left(\Vnl|\psi\rangle\right)
   - \Vnl\left(\hat{r}_\alpha|\psi\rangle\right).
  \label{eq:comm-expand}
\end{align}

\paragraph{First term: $\hat{r}_\alpha(\Vnl|\psi\rangle)$.}
Inserting the KB form of $\Vnl$:
\begin{align}
  \hat{r}_\alpha\left(\Vnl|\psi\rangle\right)
  &= \hat{r}_\alpha\sum_{I,l,m}
     |\beta^{lm}_I\rangle\, h^l_I\,
     \langle\beta^{lm}_I|\psi\rangle
  \nonumber \\
  &= \sum_{I,l,m}
     \hat{r}_\alpha|\beta^{lm}_I\rangle\, h^l_I\,
     \langle\beta^{lm}_I|\psi\rangle,
  \label{eq:term1}
\end{align}
where $\hat{r}_\alpha$ acts on the ket $|\beta^{lm}_I\rangle$
(multiplying the position wavefunction $\beta^{lm}_I(\mathbf{r})$
by $r_\alpha$), while the scalar overlap
$c_{I,lm} \equiv \langle\beta^{lm}_I|\psi\rangle
= \int\beta^{lm*}_I(\mathbf{r}')\psi(\mathbf{r}')\,\mathrm{d}^3r'$
is unaffected.

\paragraph{Second term: $\Vnl(\hat{r}_\alpha|\psi\rangle)$.}
\begin{align}
  \Vnl\left(\hat{r}_\alpha|\psi\rangle\right)
  &= \sum_{I,l,m}
     |\beta^{lm}_I\rangle\, h^l_I\,
     \langle\beta^{lm}_I|\hat{r}_\alpha|\psi\rangle.
  \label{eq:term2}
\end{align}
Since $\hat{r}_\alpha$ is Hermitian,
$\langle\beta^{lm}_I|\hat{r}_\alpha|\psi\rangle
= \langle\hat{r}_\alpha\beta^{lm}_I|\psi\rangle
= \int r_\alpha\,\beta^{lm*}_I(\mathbf{r})\,\psi(\mathbf{r})\,
\mathrm{d}^3r$.
We denote this as $\langle r_\alpha\beta^{lm}_I|\psi\rangle$,
where $r_\alpha\beta^{lm}_I$ means the function
$r_\alpha\beta^{lm*}_I(\mathbf{r})$ in position space.

\paragraph{The commutator.}
Subtracting eq.~(\ref{eq:term2}) from eq.~(\ref{eq:term1}):
\begin{align}
  [\hat{r}_\alpha,\Vnl]|\psi\rangle
  &= \sum_{I,l,m} h^l_I
     \Big[
     \hat{r}_\alpha|\beta^{lm}_I\rangle\,\langle\beta^{lm}_I|\psi\rangle
   - |\beta^{lm}_I\rangle\,\langle\hat{r}_\alpha\beta^{lm}_I|\psi\rangle
     \Big].
  \label{eq:comm-intermediate}
\end{align}
Using the shorthand $|r_\alpha\beta^{lm}_I\rangle \equiv
\hat{r}_\alpha|\beta^{lm}_I\rangle$ for the position-weighted
projector ket, this is
\begin{equation}
  \boxed{
  [\hat{r}_\alpha,\Vnl]|\psi\rangle
  = \sum_{I,l,m} h^l_I\!\left[
    |r_\alpha\beta^{lm}_I\rangle\langle\beta^{lm}_I|\psi\rangle
  - |\beta^{lm}_I\rangle\langle r_\alpha\beta^{lm}_I|\psi\rangle
    \right].
  }
  \label{eq:comm-result}
\end{equation}
This is a rank-2 update (in the language of matrix algebra): the
commutator of a separable operator with $\hat{r}_\alpha$ is
again separable, but now involves two types of projectors — one
weighted by $r_\alpha$ in the ket, one in the bra.

Equation~(\ref{eq:comm-result}) has a transparent structure:
\begin{itemize}
  \item The \emph{first term}
    $|r_\alpha\beta^{lm}_I\rangle\langle\beta^{lm}_I|\psi\rangle$
    keeps the unmodified projector in the bra and weights the ket
    by $r_\alpha$.
    Physically: the electron is projected onto the original KB
    projector at $\mathbf{r}'$ and then the result is multiplied
    by $r_\alpha$ at the field point $\mathbf{r}$.
  \item The \emph{second term}
    $|\beta^{lm}_I\rangle\langle r_\alpha\beta^{lm}_I|\psi\rangle$
    keeps the unmodified projector ket but weights the bra by
    $r_\alpha$.
    Physically: the electron is projected onto the
    $r_\alpha$-weighted projector (measuring overlap with the
    position-weighted projector) and the result propagates with
    the original projector shape.
\end{itemize}
The two terms have opposite signs and together ensure that the
commutator operator is anti-Hermitian, consistent with
$[\hat{r}_\alpha, \Vnl]^\dag = -[\hat{r}_\alpha, \Vnl]$ for
Hermitian $\hat{r}_\alpha$ and $\Vnl$.

\paragraph{Matrix element between Kohn-Sham orbitals.}
Taking the matrix element of eq.~(\ref{eq:comm-result}) between
$\langle\psi_a|$ and $|\psi_b\rangle$:
\begin{align}
  \langle\psi_a|[\hat{r}_\alpha,\Vnl]|\psi_b\rangle
  &= \sum_{I,l,m} h^l_I
     \Big[
       \underbrace{\langle\psi_a|r_\alpha\beta^{lm}_I\rangle}_{
         \displaystyle A^{lm}_{I,a}}
       \underbrace{\langle\beta^{lm}_I|\psi_b\rangle}_{
         \displaystyle B^{lm}_{I,b}}
     - \underbrace{\langle\psi_a|\beta^{lm}_I\rangle}_{
         \displaystyle \overline{B^{lm}_{I,a}}}
       \underbrace{\langle r_\alpha\beta^{lm}_I|\psi_b\rangle}_{
         \displaystyle \overline{A^{lm}_{I,b}}}
     \Big],
  \label{eq:comm-matrix-element}
\end{align}
where we introduced the notation
\begin{align}
  B^{lm}_{I,b}
  &= \langle\beta^{lm}_I|\psi_b\rangle
   = \int \beta^{lm*}_I(\mathbf{r})\,\psi_b(\mathbf{r})
     \,\mathrm{d}^3r,
  \label{eq:B-overlap}
  \\
  A^{lm}_{I,b}
  &= \langle r_\alpha\beta^{lm}_I|\psi_b\rangle
   = \int r_\alpha\,\beta^{lm*}_I(\mathbf{r})\,\psi_b(\mathbf{r})
     \,\mathrm{d}^3r.
  \label{eq:A-overlap}
\end{align}
Here $A^{lm}_{I,b}$ is the overlap of the $\alpha$-position-weighted
projector $r_\alpha\beta^{lm}_I$ with the orbital $\psi_b$, and
$B^{lm}_{I,b}$ is the standard KB projector overlap.
Both are computed as spatial integrals over the real-space grid.

Equation~(\ref{eq:comm-matrix-element}) therefore requires two types
of integrals per projector per orbital: the standard KB overlap
$B^{lm}_{I,b}$ (already available from the usual nonlocal PP force
evaluation) and the new position-weighted overlap $A^{lm}_{I,b}$.

%--------------------------------------------------------------------
\subsection*{C.\; Generalization to Multiple Projectors per Channel}
\label{app:multi-projector}
%--------------------------------------------------------------------

In the general KB form with multiple projectors per angular momentum
channel (eq.~11 of Kleinman and Bylander\cite{kleinman-bylander-1982}),
\begin{equation}
  \Vnl = \sum_I\sum_{l,m}\sum_{i,j}
  |\beta^{lm}_{I,i}\rangle\, h^l_{I,ij}\,
  \langle\beta^{lm}_{I,j}|,
  \label{eq:KB-general}
\end{equation}
where $i,j$ are radial projector indices and $h^l_{I,ij}$ is the
(generally non-diagonal) strength matrix.
The commutator result of eq.~(\ref{eq:comm-result}) extends directly
to
\begin{equation}
  [\hat{r}_\alpha,\Vnl]|\psi\rangle
  = \sum_{I,l,m}\sum_{i,j} h^l_{I,ij}\!\left[
    |r_\alpha\beta^{lm}_{I,i}\rangle
    \langle\beta^{lm}_{I,j}|\psi\rangle
  - |\beta^{lm}_{I,i}\rangle
    \langle r_\alpha\beta^{lm}_{I,j}|\psi\rangle
    \right].
  \label{eq:comm-result-general}
\end{equation}
The algebraic structure is identical; only the matrix of strength
coefficients $h^l_{I,ij}$ becomes non-diagonal in the projector
indices.
For the norm-conserving pseudopotentials used in \texttt{RMG}, the
single-projector case (eq.~\ref{eq:comm-result}) applies.

%--------------------------------------------------------------------
\subsection*{D.\; Nonlocal Contribution to the Current Density}
\label{app:Jnl-expression}
%--------------------------------------------------------------------

The nonlocal current density contribution to
eq.~(\ref{eq:J-total}) is defined through
eq.~(\ref{eq:Jnl-continuity}).
We now verify that this definition is consistent with the velocity
operator result by showing that the volume integral of
$\mathbf{J}_\mathrm{nl}$ gives the expected macroscopic contribution.

Using the single-particle density matrix $P(\mathbf{r},\mathbf{r}')
= \sum_n f_n \psi_n(\mathbf{r})\psi_n^*(\mathbf{r}')$, the
nonlocal contribution to the macroscopic current along direction
$\alpha$ is
\begin{align}
  J_{\mathrm{nl},\alpha}(t)
  &= \int_\Omega\!
     \operatorname{Re}\!\left[
       \psi^*(\mathbf{r},t)\,
       \frac{1}{\imath}[\hat{r}_\alpha,\Vnl]\,
       \psi(\mathbf{r},t)
     \right]\mathrm{d}^3r
  \nonumber \\
  &= \frac{1}{\imath}\sum_{I,l,m} h^l_I\,
     \operatorname{Re}\!\Bigg[
       \int\!\psi^*(\mathbf{r})\,r_\alpha\,\beta^{lm}_I(\mathbf{r})
       \,\mathrm{d}^3r
       \cdot
       \int\!\beta^{lm*}_I(\mathbf{r}')\,\psi(\mathbf{r}')
       \,\mathrm{d}^3r'
  \nonumber \\
  &\qquad\qquad\qquad\quad
   - \int\!\psi^*(\mathbf{r})\,\beta^{lm}_I(\mathbf{r})
     \,\mathrm{d}^3r
     \cdot
     \int\!r_\alpha\,\beta^{lm*}_I(\mathbf{r}')\,\psi(\mathbf{r}')
     \,\mathrm{d}^3r'
     \Bigg].
  \label{eq:Jnl-full}
\end{align}
In the density matrix notation of Section~\ref{sec:current-dm-sec},
this is evaluated as the matrix element
$\langle\psi_a|\frac{1}{\imath}[\hat{r}_\alpha,\Vnl]|\psi_b\rangle$
from eq.~(\ref{eq:comm-matrix-element}), contracted with $P_{ab}(t)$.
No additional real-space integration is required at each time step
beyond what is already computed in building $\Jmat{\alpha}_{ab}$.

%--------------------------------------------------------------------
\subsection*{E.\; Connection to \texttt{AppNls\_0xyz()} in RMG}
\label{app:code-connection}
%--------------------------------------------------------------------

In \texttt{RMG}, the commutator $[\hat{r}_\alpha,\Vnl]|\psi_b\rangle$
is evaluated by the function \texttt{AppNls\_0xyz()}, called from
\texttt{CurrentNlpp()} for directions $\alpha \in \{x,y,z\}$.

The function computes exactly eq.~(\ref{eq:comm-result}) on the
real-space grid.
For each orbital $\psi_b$ and each atom $I$, it evaluates:
\begin{enumerate}
  \item The standard projector overlap
    $B^{lm}_{I,b} = \langle\beta^{lm}_I|\psi_b\rangle$
    (a sum over grid points in the support of $\beta^{lm}_I$);
  \item The position-weighted overlap
    $A^{lm}_{I,b} = \langle r_\alpha\beta^{lm}_I|\psi_b\rangle$
    (the same sum but with each grid point contribution weighted
    by its $r_\alpha$ coordinate relative to atom $I$);
  \item The output array as
    $[\hat{r}_\alpha,\Vnl]\psi_b(\mathbf{r})
    = \sum_{I,l,m} h^l_I\!\left[
      r_\alpha\beta^{lm}_I(\mathbf{r})\cdot B^{lm}_{I,b}
      - \beta^{lm}_I(\mathbf{r})\cdot A^{lm}_{I,b}
    \right]$,
    stored in the work array \texttt{nv[]}.
\end{enumerate}
The result is passed to \texttt{CurrentNlpp()}, which computes the
matrix element $\langle\psi_a|\texttt{nv}\rangle$ via a BLAS
\texttt{gemm} operation, multiplies by \texttt{I\_t} $= \imath$,
and adds the result to \texttt{Palpha\_matrix}.
As shown in Section~\ref{sec:imath-convention}, the combined factor
$\imath \times \frac{1}{\imath} = 1$ ensures that the nonlocal PP
correction enters the total current matrix
$\boldsymbol{\mathcal{J}}^\alpha$ with unit weight, consistent with
eq.~(\ref{eq:Jmat-full}).

A self-consistency check: at the ground state, when all orbitals
$\psi_n$ are real-valued eigenstates, the density matrix is
$P_{ab} = f_a\delta_{ab}$ and the imaginary part of
$\langle\psi_a|\frac{1}{\imath}[\hat{r}_\alpha,\Vnl]|\psi_a\rangle$
should vanish because $[\hat{r}_\alpha,\Vnl]$ is anti-Hermitian
(the diagonal elements of an anti-Hermitian operator are purely
imaginary for complex orbitals and zero for real orbitals).
Thus $J_\mathrm{nl}(t=0) = 0$, consistent with vanishing ground-state
current.
This provides a simple numerical test that can be run on the
ground-state density matrix before starting the RT-TDDFT propagation.
 goes in the appendix section of paper2_main.tex.
%%
%% Structure:
%%   A. Gauge transformation V^nl under LG -> VG: derives the coupling
%%      -i A_alpha [r_alpha, V^nl] that appears in H^VG
%%   B. Explicit algebra of [r_alpha, V^nl] for single-projector KB PP
%%   C. Generalization to multiple projectors per channel
%%   D. Nonlocal current density J_nl from the commutator
%%   E. Connection to AppNls_0xyz() in RMG
%%
%% CONVENTIONS: same as sec_theory_current.tex (atomic units, e_electron=-1)
%% KEY SOURCE: NEAT-project-notes_main.tex lines 785-851 (Starace derivation)
%%             Kleinman-Bylander (1982); Mattiat-Luber (2022) eq.12
%%%%%%%%%%%%%%%%%%%%%%%%%%%%%%%%%%%%%%%%%%%%%%%%%%%%%%%%%%%%%%%%%%%%%%%%%%%%%%%

\section{Gauge Transformation of the Nonlocal Pseudopotential
  and the Commutator $[\hat{r}_\alpha, \Vnl]$}
\label{app:gauge_nlpp}

This appendix derives two connected results:
(A)~the transformation of the nonlocal Kleinman-Bylander (KB)
pseudopotential from the length gauge to the velocity gauge,
establishing why the commutator $[\hat{r}_\alpha, \Vnl]$ enters the
velocity gauge Hamiltonian; and (B)~the explicit algebraic evaluation
of the commutator for a separable pseudopotential, providing the
formula implemented in \texttt{AppNls\_0xyz()} in \texttt{RMG}.

%--------------------------------------------------------------------
\subsection*{A.\; Gauge Transformation of $\Vnl$ from Length to
  Velocity Gauge}
\label{app:gauge-transform-Vnl}
%--------------------------------------------------------------------

Under the gauge transformation from length gauge (LG) to velocity
gauge (VG), the wavefunction and all operators transform by the
unitary operator\cite{Luber-rt-tddft,Mattiat-Luber-2022}
\begin{equation}
  U(\mathbf{r},t)
  = \exp\!\left(\imath\frac{q}{\hbar c}\chi(\mathbf{r},t)\right)
  \stackrel{\mathrm{a.u.}}{=}
  \exp(-\imath\chi(\mathbf{r},t)),
  \label{eq:U-gauge}
\end{equation}
where in atomic units with $\hbar = c = 1$ and $q = -1$ for
electrons, $q/\hbar c = -1$.
The gauge function for the LG $\to$ VG transformation is
\begin{equation}
  \chi(\mathbf{r},t) = -\mathbf{r}\cdot\Avec(t),
  \label{eq:chi-vg}
\end{equation}
where $\Avec(t) = -\int_0^t \mathbf{E}(t')\,\mathrm{d}t'$ is the
vector potential chosen so that $\phi'=0$ in the velocity
gauge.\cite{Luber-rt-tddft,rt-DFTB-Wong-2023}
With $\chi = -\mathbf{r}\cdot\Avec$, the transformation operator is
\begin{equation}
  U(\mathbf{r},t)
  = e^{-\imath(-\mathbf{r}\cdot\Avec)}
  = e^{\imath\mathbf{r}\cdot\Avec(t)}.
  \label{eq:U-explicit}
\end{equation}

Any operator $\hat{O}$ transforms as $\hat{O}' = U\hat{O}U^\dag$.
For a \emph{local} operator $f(\mathbf{r})$, the transformation is
trivial: $U f(\mathbf{r})U^\dag = f(\mathbf{r})$, because
$U(\mathbf{r})$ commutes with any function of $\mathbf{r}$.

For the \emph{nonlocal} pseudopotential $\Vnl$, the two operators
$U$ and $\Vnl$ do not commute, and the transformed potential is
\begin{equation}
  \tilde{V}^{\mathrm{nl}}
  = U\Vnl U^\dag
  = e^{\imath\mathbf{r}\cdot\Avec}\Vnl
    e^{-\imath\mathbf{r}\cdot\Avec}.
  \label{eq:Vnl-transformed}
\end{equation}
To extract the first-order coupling to $\Avec$, we use the
Baker-Campbell-Hausdorff (BCH) expansion truncated at first order
in $\Avec$ (valid within linear response):
\begin{align}
  e^{\imath\mathbf{r}\cdot\Avec}\Vnl e^{-\imath\mathbf{r}\cdot\Avec}
  &= \Vnl
   + \imath[\mathbf{r}\cdot\Avec,\, \Vnl]
   + \mathcal{O}(\Avec^2)
  \nonumber \\
  &= \Vnl
   + \imath A_\alpha(t)[\hat{r}_\alpha,\Vnl]
   + \mathcal{O}(\Avec^2),
  \label{eq:Vnl-BCH}
\end{align}
where repeated Cartesian index $\alpha$ is summed.
The velocity gauge Hamiltonian therefore differs from the length
gauge Hamiltonian by the first-order correction
\begin{equation}
  H^{\mathrm{VG}} - H^{\mathrm{LG}}
  \supset \imath A_\alpha(t)[\hat{r}_\alpha,\Vnl]
  + \mathcal{O}(\Avec^2).
  \label{eq:H-VG-correction}
\end{equation}
This correction is the origin of the nonlocal pseudopotential term
in the velocity gauge Hamiltonian.
When combined with the paramagnetic coupling
$-\Avec\cdot\hat{\mathbf{p}}$ (from the kinetic energy), the
total field-electron coupling in the velocity gauge is
\begin{equation}
  V_{\mathrm{field}}^{\mathrm{VG}}
  = -A_\alpha(t)\hat{p}_\alpha
    + \imath A_\alpha(t)[\hat{r}_\alpha,\Vnl]
  = A_\alpha(t)\!\left(-\hat{p}_\alpha
    + \imath[\hat{r}_\alpha,\Vnl]\right),
  \label{eq:V-field-VG}
\end{equation}
which can be written as $-A_\alpha(t)\hat{v}_\alpha$ upon
identifying $\hat{v}_\alpha = \hat{p}_\alpha -
\imath[\hat{r}_\alpha,\Vnl] = \hat{p}_\alpha +
\frac{1}{\imath}[\hat{r}_\alpha,\Vnl]\cdot(-1)$.
(The sign is consistent with $\hat{v}_\alpha = \hat{p}_\alpha +
\frac{1}{\imath}[\hat{r}_\alpha,\Vnl]$ at $A=0$, as derived in
Section~\ref{sec:velocity-eom}.)

%--------------------------------------------------------------------
\subsection*{B.\; Explicit Evaluation of $[\hat{r}_\alpha, \Vnl]$
  for Kleinman-Bylander Pseudopotentials}
\label{app:KB-commutator}
%--------------------------------------------------------------------

The Kleinman-Bylander separable pseudopotential for atom $I$ in the
single-projector-per-channel approximation
is\cite{kleinman-bylander-1982}
\begin{equation}
  \Vnl = \sum_I\sum_{l,m}
  |\beta^{lm}_I\rangle\, h^l_I\,
  \langle\beta^{lm}_I|,
  \label{eq:KB-def}
\end{equation}
where $|\beta^{lm}_I\rangle$ is the KB projector for atom $I$
with angular momentum quantum numbers $(l,m)$, centered at
$\mathbf{R}_I$, and $h^l_I$ is the (real, scalar) strength
coefficient for angular momentum channel $l$.
In position representation, the projector wavefunction is
$\beta^{lm}_I(\mathbf{r}) = \langle\mathbf{r}|\beta^{lm}_I\rangle$.
The sum over $(l,m)$ runs over all angular momentum channels and
magnetic quantum numbers included in the pseudopotential.

We evaluate the commutator $[\hat{r}_\alpha, \Vnl]$ by direct
computation, acting on an arbitrary state $|\psi\rangle$:
\begin{align}
  [\hat{r}_\alpha, \Vnl]|\psi\rangle
  &= \hat{r}_\alpha\left(\Vnl|\psi\rangle\right)
   - \Vnl\left(\hat{r}_\alpha|\psi\rangle\right).
  \label{eq:comm-expand}
\end{align}

\paragraph{First term: $\hat{r}_\alpha(\Vnl|\psi\rangle)$.}
Inserting the KB form of $\Vnl$:
\begin{align}
  \hat{r}_\alpha\left(\Vnl|\psi\rangle\right)
  &= \hat{r}_\alpha\sum_{I,l,m}
     |\beta^{lm}_I\rangle\, h^l_I\,
     \langle\beta^{lm}_I|\psi\rangle
  \nonumber \\
  &= \sum_{I,l,m}
     \hat{r}_\alpha|\beta^{lm}_I\rangle\, h^l_I\,
     \langle\beta^{lm}_I|\psi\rangle,
  \label{eq:term1}
\end{align}
where $\hat{r}_\alpha$ acts on the ket $|\beta^{lm}_I\rangle$
(multiplying the position wavefunction $\beta^{lm}_I(\mathbf{r})$
by $r_\alpha$), while the scalar overlap
$c_{I,lm} \equiv \langle\beta^{lm}_I|\psi\rangle
= \int\beta^{lm*}_I(\mathbf{r}')\psi(\mathbf{r}')\,\mathrm{d}^3r'$
is unaffected.

\paragraph{Second term: $\Vnl(\hat{r}_\alpha|\psi\rangle)$.}
\begin{align}
  \Vnl\left(\hat{r}_\alpha|\psi\rangle\right)
  &= \sum_{I,l,m}
     |\beta^{lm}_I\rangle\, h^l_I\,
     \langle\beta^{lm}_I|\hat{r}_\alpha|\psi\rangle.
  \label{eq:term2}
\end{align}
Since $\hat{r}_\alpha$ is Hermitian,
$\langle\beta^{lm}_I|\hat{r}_\alpha|\psi\rangle
= \langle\hat{r}_\alpha\beta^{lm}_I|\psi\rangle
= \int r_\alpha\,\beta^{lm*}_I(\mathbf{r})\,\psi(\mathbf{r})\,
\mathrm{d}^3r$.
We denote this as $\langle r_\alpha\beta^{lm}_I|\psi\rangle$,
where $r_\alpha\beta^{lm}_I$ means the function
$r_\alpha\beta^{lm*}_I(\mathbf{r})$ in position space.

\paragraph{The commutator.}
Subtracting eq.~(\ref{eq:term2}) from eq.~(\ref{eq:term1}):
\begin{align}
  [\hat{r}_\alpha,\Vnl]|\psi\rangle
  &= \sum_{I,l,m} h^l_I
     \Big[
     \hat{r}_\alpha|\beta^{lm}_I\rangle\,\langle\beta^{lm}_I|\psi\rangle
   - |\beta^{lm}_I\rangle\,\langle\hat{r}_\alpha\beta^{lm}_I|\psi\rangle
     \Big].
  \label{eq:comm-intermediate}
\end{align}
Using the shorthand $|r_\alpha\beta^{lm}_I\rangle \equiv
\hat{r}_\alpha|\beta^{lm}_I\rangle$ for the position-weighted
projector ket, this is
\begin{equation}
  \boxed{
  [\hat{r}_\alpha,\Vnl]|\psi\rangle
  = \sum_{I,l,m} h^l_I\!\left[
    |r_\alpha\beta^{lm}_I\rangle\langle\beta^{lm}_I|\psi\rangle
  - |\beta^{lm}_I\rangle\langle r_\alpha\beta^{lm}_I|\psi\rangle
    \right].
  }
  \label{eq:comm-result}
\end{equation}
This is a rank-2 update (in the language of matrix algebra): the
commutator of a separable operator with $\hat{r}_\alpha$ is
again separable, but now involves two types of projectors — one
weighted by $r_\alpha$ in the ket, one in the bra.

Equation~(\ref{eq:comm-result}) has a transparent structure:
\begin{itemize}
  \item The \emph{first term}
    $|r_\alpha\beta^{lm}_I\rangle\langle\beta^{lm}_I|\psi\rangle$
    keeps the unmodified projector in the bra and weights the ket
    by $r_\alpha$.
    Physically: the electron is projected onto the original KB
    projector at $\mathbf{r}'$ and then the result is multiplied
    by $r_\alpha$ at the field point $\mathbf{r}$.
  \item The \emph{second term}
    $|\beta^{lm}_I\rangle\langle r_\alpha\beta^{lm}_I|\psi\rangle$
    keeps the unmodified projector ket but weights the bra by
    $r_\alpha$.
    Physically: the electron is projected onto the
    $r_\alpha$-weighted projector (measuring overlap with the
    position-weighted projector) and the result propagates with
    the original projector shape.
\end{itemize}
The two terms have opposite signs and together ensure that the
commutator operator is anti-Hermitian, consistent with
$[\hat{r}_\alpha, \Vnl]^\dag = -[\hat{r}_\alpha, \Vnl]$ for
Hermitian $\hat{r}_\alpha$ and $\Vnl$.

\paragraph{Matrix element between Kohn-Sham orbitals.}
Taking the matrix element of eq.~(\ref{eq:comm-result}) between
$\langle\psi_a|$ and $|\psi_b\rangle$:
\begin{align}
  \langle\psi_a|[\hat{r}_\alpha,\Vnl]|\psi_b\rangle
  &= \sum_{I,l,m} h^l_I
     \Big[
       \underbrace{\langle\psi_a|r_\alpha\beta^{lm}_I\rangle}_{
         \displaystyle A^{lm}_{I,a}}
       \underbrace{\langle\beta^{lm}_I|\psi_b\rangle}_{
         \displaystyle B^{lm}_{I,b}}
     - \underbrace{\langle\psi_a|\beta^{lm}_I\rangle}_{
         \displaystyle \overline{B^{lm}_{I,a}}}
       \underbrace{\langle r_\alpha\beta^{lm}_I|\psi_b\rangle}_{
         \displaystyle \overline{A^{lm}_{I,b}}}
     \Big],
  \label{eq:comm-matrix-element}
\end{align}
where we introduced the notation
\begin{align}
  B^{lm}_{I,b}
  &= \langle\beta^{lm}_I|\psi_b\rangle
   = \int \beta^{lm*}_I(\mathbf{r})\,\psi_b(\mathbf{r})
     \,\mathrm{d}^3r,
  \label{eq:B-overlap}
  \\
  A^{lm}_{I,b}
  &= \langle r_\alpha\beta^{lm}_I|\psi_b\rangle
   = \int r_\alpha\,\beta^{lm*}_I(\mathbf{r})\,\psi_b(\mathbf{r})
     \,\mathrm{d}^3r.
  \label{eq:A-overlap}
\end{align}
Here $A^{lm}_{I,b}$ is the overlap of the $\alpha$-position-weighted
projector $r_\alpha\beta^{lm}_I$ with the orbital $\psi_b$, and
$B^{lm}_{I,b}$ is the standard KB projector overlap.
Both are computed as spatial integrals over the real-space grid.

Equation~(\ref{eq:comm-matrix-element}) therefore requires two types
of integrals per projector per orbital: the standard KB overlap
$B^{lm}_{I,b}$ (already available from the usual nonlocal PP force
evaluation) and the new position-weighted overlap $A^{lm}_{I,b}$.

%--------------------------------------------------------------------
\subsection*{C.\; Generalization to Multiple Projectors per Channel}
\label{app:multi-projector}
%--------------------------------------------------------------------

In the general KB form with multiple projectors per angular momentum
channel (eq.~11 of Kleinman and Bylander\cite{kleinman-bylander-1982}),
\begin{equation}
  \Vnl = \sum_I\sum_{l,m}\sum_{i,j}
  |\beta^{lm}_{I,i}\rangle\, h^l_{I,ij}\,
  \langle\beta^{lm}_{I,j}|,
  \label{eq:KB-general}
\end{equation}
where $i,j$ are radial projector indices and $h^l_{I,ij}$ is the
(generally non-diagonal) strength matrix.
The commutator result of eq.~(\ref{eq:comm-result}) extends directly
to
\begin{equation}
  [\hat{r}_\alpha,\Vnl]|\psi\rangle
  = \sum_{I,l,m}\sum_{i,j} h^l_{I,ij}\!\left[
    |r_\alpha\beta^{lm}_{I,i}\rangle
    \langle\beta^{lm}_{I,j}|\psi\rangle
  - |\beta^{lm}_{I,i}\rangle
    \langle r_\alpha\beta^{lm}_{I,j}|\psi\rangle
    \right].
  \label{eq:comm-result-general}
\end{equation}
The algebraic structure is identical; only the matrix of strength
coefficients $h^l_{I,ij}$ becomes non-diagonal in the projector
indices.
For the norm-conserving pseudopotentials used in \texttt{RMG}, the
single-projector case (eq.~\ref{eq:comm-result}) applies.

%--------------------------------------------------------------------
\subsection*{D.\; Nonlocal Contribution to the Current Density}
\label{app:Jnl-expression}
%--------------------------------------------------------------------

The nonlocal current density contribution to
eq.~(\ref{eq:J-total}) is defined through
eq.~(\ref{eq:Jnl-continuity}).
We now verify that this definition is consistent with the velocity
operator result by showing that the volume integral of
$\mathbf{J}_\mathrm{nl}$ gives the expected macroscopic contribution.

Using the single-particle density matrix $P(\mathbf{r},\mathbf{r}')
= \sum_n f_n \psi_n(\mathbf{r})\psi_n^*(\mathbf{r}')$, the
nonlocal contribution to the macroscopic current along direction
$\alpha$ is
\begin{align}
  J_{\mathrm{nl},\alpha}(t)
  &= \int_\Omega\!
     \operatorname{Re}\!\left[
       \psi^*(\mathbf{r},t)\,
       \frac{1}{\imath}[\hat{r}_\alpha,\Vnl]\,
       \psi(\mathbf{r},t)
     \right]\mathrm{d}^3r
  \nonumber \\
  &= \frac{1}{\imath}\sum_{I,l,m} h^l_I\,
     \operatorname{Re}\!\Bigg[
       \int\!\psi^*(\mathbf{r})\,r_\alpha\,\beta^{lm}_I(\mathbf{r})
       \,\mathrm{d}^3r
       \cdot
       \int\!\beta^{lm*}_I(\mathbf{r}')\,\psi(\mathbf{r}')
       \,\mathrm{d}^3r'
  \nonumber \\
  &\qquad\qquad\qquad\quad
   - \int\!\psi^*(\mathbf{r})\,\beta^{lm}_I(\mathbf{r})
     \,\mathrm{d}^3r
     \cdot
     \int\!r_\alpha\,\beta^{lm*}_I(\mathbf{r}')\,\psi(\mathbf{r}')
     \,\mathrm{d}^3r'
     \Bigg].
  \label{eq:Jnl-full}
\end{align}
In the density matrix notation of Section~\ref{sec:current-dm-sec},
this is evaluated as the matrix element
$\langle\psi_a|\frac{1}{\imath}[\hat{r}_\alpha,\Vnl]|\psi_b\rangle$
from eq.~(\ref{eq:comm-matrix-element}), contracted with $P_{ab}(t)$.
No additional real-space integration is required at each time step
beyond what is already computed in building $\Jmat{\alpha}_{ab}$.

%--------------------------------------------------------------------
\subsection*{E.\; Connection to \texttt{AppNls\_0xyz()} in RMG}
\label{app:code-connection}
%--------------------------------------------------------------------

In \texttt{RMG}, the commutator $[\hat{r}_\alpha,\Vnl]|\psi_b\rangle$
is evaluated by the function \texttt{AppNls\_0xyz()}, called from
\texttt{CurrentNlpp()} for directions $\alpha \in \{x,y,z\}$.

The function computes exactly eq.~(\ref{eq:comm-result}) on the
real-space grid.
For each orbital $\psi_b$ and each atom $I$, it evaluates:
\begin{enumerate}
  \item The standard projector overlap
    $B^{lm}_{I,b} = \langle\beta^{lm}_I|\psi_b\rangle$
    (a sum over grid points in the support of $\beta^{lm}_I$);
  \item The position-weighted overlap
    $A^{lm}_{I,b} = \langle r_\alpha\beta^{lm}_I|\psi_b\rangle$
    (the same sum but with each grid point contribution weighted
    by its $r_\alpha$ coordinate relative to atom $I$);
  \item The output array as
    $[\hat{r}_\alpha,\Vnl]\psi_b(\mathbf{r})
    = \sum_{I,l,m} h^l_I\!\left[
      r_\alpha\beta^{lm}_I(\mathbf{r})\cdot B^{lm}_{I,b}
      - \beta^{lm}_I(\mathbf{r})\cdot A^{lm}_{I,b}
    \right]$,
    stored in the work array \texttt{nv[]}.
\end{enumerate}
The result is passed to \texttt{CurrentNlpp()}, which computes the
matrix element $\langle\psi_a|\texttt{nv}\rangle$ via a BLAS
\texttt{gemm} operation, multiplies by \texttt{I\_t} $= \imath$,
and adds the result to \texttt{Palpha\_matrix}.
As shown in Section~\ref{sec:imath-convention}, the combined factor
$\imath \times \frac{1}{\imath} = 1$ ensures that the nonlocal PP
correction enters the total current matrix
$\boldsymbol{\mathcal{J}}^\alpha$ with unit weight, consistent with
eq.~(\ref{eq:Jmat-full}).

A self-consistency check: at the ground state, when all orbitals
$\psi_n$ are real-valued eigenstates, the density matrix is
$P_{ab} = f_a\delta_{ab}$ and the imaginary part of
$\langle\psi_a|\frac{1}{\imath}[\hat{r}_\alpha,\Vnl]|\psi_a\rangle$
should vanish because $[\hat{r}_\alpha,\Vnl]$ is anti-Hermitian
(the diagonal elements of an anti-Hermitian operator are purely
imaginary for complex orbitals and zero for real orbitals).
Thus $J_\mathrm{nl}(t=0) = 0$, consistent with vanishing ground-state
current.
This provides a simple numerical test that can be run on the
ground-state density matrix before starting the RT-TDDFT propagation.
 goes in the appendix section of paper2_main.tex.
%%
%% Structure:
%%   A. Gauge transformation V^nl under LG -> VG: derives the coupling
%%      -i A_alpha [r_alpha, V^nl] that appears in H^VG
%%   B. Explicit algebra of [r_alpha, V^nl] for single-projector KB PP
%%   C. Generalization to multiple projectors per channel
%%   D. Nonlocal current density J_nl from the commutator
%%   E. Connection to AppNls_0xyz() in RMG
%%
%% CONVENTIONS: same as sec_theory_current.tex (atomic units, e_electron=-1)
%% KEY SOURCE: NEAT-project-notes_main.tex lines 785-851 (Starace derivation)
%%             Kleinman-Bylander (1982); Mattiat-Luber (2022) eq.12
%%%%%%%%%%%%%%%%%%%%%%%%%%%%%%%%%%%%%%%%%%%%%%%%%%%%%%%%%%%%%%%%%%%%%%%%%%%%%%%

\section{Gauge Transformation of the Nonlocal Pseudopotential
  and the Commutator $[\hat{r}_\alpha, \Vnl]$}
\label{app:gauge_nlpp}

This appendix derives two connected results:
(A)~the transformation of the nonlocal Kleinman-Bylander (KB)
pseudopotential from the length gauge to the velocity gauge,
establishing why the commutator $[\hat{r}_\alpha, \Vnl]$ enters the
velocity gauge Hamiltonian; and (B)~the explicit algebraic evaluation
of the commutator for a separable pseudopotential, providing the
formula implemented in \texttt{AppNls\_0xyz()} in \texttt{RMG}.

%--------------------------------------------------------------------
\subsection*{A.\; Gauge Transformation of $\Vnl$ from Length to
  Velocity Gauge}
\label{app:gauge-transform-Vnl}
%--------------------------------------------------------------------

Under the gauge transformation from length gauge (LG) to velocity
gauge (VG), the wavefunction and all operators transform by the
unitary operator\cite{Luber-rt-tddft,Mattiat-Luber-2022}
\begin{equation}
  U(\mathbf{r},t)
  = \exp\!\left(\imath\frac{q}{\hbar c}\chi(\mathbf{r},t)\right)
  \stackrel{\mathrm{a.u.}}{=}
  \exp(-\imath\chi(\mathbf{r},t)),
  \label{eq:U-gauge}
\end{equation}
where in atomic units with $\hbar = c = 1$ and $q = -1$ for
electrons, $q/\hbar c = -1$.
The gauge function for the LG $\to$ VG transformation is
\begin{equation}
  \chi(\mathbf{r},t) = -\mathbf{r}\cdot\Avec(t),
  \label{eq:chi-vg}
\end{equation}
where $\Avec(t) = -\int_0^t \mathbf{E}(t')\,\mathrm{d}t'$ is the
vector potential chosen so that $\phi'=0$ in the velocity
gauge.\cite{Luber-rt-tddft,rt-DFTB-Wong-2023}
With $\chi = -\mathbf{r}\cdot\Avec$, the transformation operator is
\begin{equation}
  U(\mathbf{r},t)
  = e^{-\imath(-\mathbf{r}\cdot\Avec)}
  = e^{\imath\mathbf{r}\cdot\Avec(t)}.
  \label{eq:U-explicit}
\end{equation}

Any operator $\hat{O}$ transforms as $\hat{O}' = U\hat{O}U^\dag$.
For a \emph{local} operator $f(\mathbf{r})$, the transformation is
trivial: $U f(\mathbf{r})U^\dag = f(\mathbf{r})$, because
$U(\mathbf{r})$ commutes with any function of $\mathbf{r}$.

For the \emph{nonlocal} pseudopotential $\Vnl$, the two operators
$U$ and $\Vnl$ do not commute, and the transformed potential is
\begin{equation}
  \tilde{V}^{\mathrm{nl}}
  = U\Vnl U^\dag
  = e^{\imath\mathbf{r}\cdot\Avec}\Vnl
    e^{-\imath\mathbf{r}\cdot\Avec}.
  \label{eq:Vnl-transformed}
\end{equation}
To extract the first-order coupling to $\Avec$, we use the
Baker-Campbell-Hausdorff (BCH) expansion truncated at first order
in $\Avec$ (valid within linear response):
\begin{align}
  e^{\imath\mathbf{r}\cdot\Avec}\Vnl e^{-\imath\mathbf{r}\cdot\Avec}
  &= \Vnl
   + \imath[\mathbf{r}\cdot\Avec,\, \Vnl]
   + \mathcal{O}(\Avec^2)
  \nonumber \\
  &= \Vnl
   + \imath A_\alpha(t)[\hat{r}_\alpha,\Vnl]
   + \mathcal{O}(\Avec^2),
  \label{eq:Vnl-BCH}
\end{align}
where repeated Cartesian index $\alpha$ is summed.
The velocity gauge Hamiltonian therefore differs from the length
gauge Hamiltonian by the first-order correction
\begin{equation}
  H^{\mathrm{VG}} - H^{\mathrm{LG}}
  \supset \imath A_\alpha(t)[\hat{r}_\alpha,\Vnl]
  + \mathcal{O}(\Avec^2).
  \label{eq:H-VG-correction}
\end{equation}
This correction is the origin of the nonlocal pseudopotential term
in the velocity gauge Hamiltonian.
When combined with the paramagnetic coupling
$-\Avec\cdot\hat{\mathbf{p}}$ (from the kinetic energy), the
total field-electron coupling in the velocity gauge is
\begin{equation}
  V_{\mathrm{field}}^{\mathrm{VG}}
  = -A_\alpha(t)\hat{p}_\alpha
    + \imath A_\alpha(t)[\hat{r}_\alpha,\Vnl]
  = A_\alpha(t)\!\left(-\hat{p}_\alpha
    + \imath[\hat{r}_\alpha,\Vnl]\right),
  \label{eq:V-field-VG}
\end{equation}
which can be written as $-A_\alpha(t)\hat{v}_\alpha$ upon
identifying $\hat{v}_\alpha = \hat{p}_\alpha -
\imath[\hat{r}_\alpha,\Vnl] = \hat{p}_\alpha +
\frac{1}{\imath}[\hat{r}_\alpha,\Vnl]\cdot(-1)$.
(The sign is consistent with $\hat{v}_\alpha = \hat{p}_\alpha +
\frac{1}{\imath}[\hat{r}_\alpha,\Vnl]$ at $A=0$, as derived in
Section~\ref{sec:velocity-eom}.)

%--------------------------------------------------------------------
\subsection*{B.\; Explicit Evaluation of $[\hat{r}_\alpha, \Vnl]$
  for Kleinman-Bylander Pseudopotentials}
\label{app:KB-commutator}
%--------------------------------------------------------------------

The Kleinman-Bylander separable pseudopotential for atom $I$ in the
single-projector-per-channel approximation
is\cite{kleinman-bylander-1982}
\begin{equation}
  \Vnl = \sum_I\sum_{l,m}
  |\beta^{lm}_I\rangle\, h^l_I\,
  \langle\beta^{lm}_I|,
  \label{eq:KB-def}
\end{equation}
where $|\beta^{lm}_I\rangle$ is the KB projector for atom $I$
with angular momentum quantum numbers $(l,m)$, centered at
$\mathbf{R}_I$, and $h^l_I$ is the (real, scalar) strength
coefficient for angular momentum channel $l$.
In position representation, the projector wavefunction is
$\beta^{lm}_I(\mathbf{r}) = \langle\mathbf{r}|\beta^{lm}_I\rangle$.
The sum over $(l,m)$ runs over all angular momentum channels and
magnetic quantum numbers included in the pseudopotential.

We evaluate the commutator $[\hat{r}_\alpha, \Vnl]$ by direct
computation, acting on an arbitrary state $|\psi\rangle$:
\begin{align}
  [\hat{r}_\alpha, \Vnl]|\psi\rangle
  &= \hat{r}_\alpha\left(\Vnl|\psi\rangle\right)
   - \Vnl\left(\hat{r}_\alpha|\psi\rangle\right).
  \label{eq:comm-expand}
\end{align}

\paragraph{First term: $\hat{r}_\alpha(\Vnl|\psi\rangle)$.}
Inserting the KB form of $\Vnl$:
\begin{align}
  \hat{r}_\alpha\left(\Vnl|\psi\rangle\right)
  &= \hat{r}_\alpha\sum_{I,l,m}
     |\beta^{lm}_I\rangle\, h^l_I\,
     \langle\beta^{lm}_I|\psi\rangle
  \nonumber \\
  &= \sum_{I,l,m}
     \hat{r}_\alpha|\beta^{lm}_I\rangle\, h^l_I\,
     \langle\beta^{lm}_I|\psi\rangle,
  \label{eq:term1}
\end{align}
where $\hat{r}_\alpha$ acts on the ket $|\beta^{lm}_I\rangle$
(multiplying the position wavefunction $\beta^{lm}_I(\mathbf{r})$
by $r_\alpha$), while the scalar overlap
$c_{I,lm} \equiv \langle\beta^{lm}_I|\psi\rangle
= \int\beta^{lm*}_I(\mathbf{r}')\psi(\mathbf{r}')\,\mathrm{d}^3r'$
is unaffected.

\paragraph{Second term: $\Vnl(\hat{r}_\alpha|\psi\rangle)$.}
\begin{align}
  \Vnl\left(\hat{r}_\alpha|\psi\rangle\right)
  &= \sum_{I,l,m}
     |\beta^{lm}_I\rangle\, h^l_I\,
     \langle\beta^{lm}_I|\hat{r}_\alpha|\psi\rangle.
  \label{eq:term2}
\end{align}
Since $\hat{r}_\alpha$ is Hermitian,
$\langle\beta^{lm}_I|\hat{r}_\alpha|\psi\rangle
= \langle\hat{r}_\alpha\beta^{lm}_I|\psi\rangle
= \int r_\alpha\,\beta^{lm*}_I(\mathbf{r})\,\psi(\mathbf{r})\,
\mathrm{d}^3r$.
We denote this as $\langle r_\alpha\beta^{lm}_I|\psi\rangle$,
where $r_\alpha\beta^{lm}_I$ means the function
$r_\alpha\beta^{lm*}_I(\mathbf{r})$ in position space.

\paragraph{The commutator.}
Subtracting eq.~(\ref{eq:term2}) from eq.~(\ref{eq:term1}):
\begin{align}
  [\hat{r}_\alpha,\Vnl]|\psi\rangle
  &= \sum_{I,l,m} h^l_I
     \Big[
     \hat{r}_\alpha|\beta^{lm}_I\rangle\,\langle\beta^{lm}_I|\psi\rangle
   - |\beta^{lm}_I\rangle\,\langle\hat{r}_\alpha\beta^{lm}_I|\psi\rangle
     \Big].
  \label{eq:comm-intermediate}
\end{align}
Using the shorthand $|r_\alpha\beta^{lm}_I\rangle \equiv
\hat{r}_\alpha|\beta^{lm}_I\rangle$ for the position-weighted
projector ket, this is
\begin{equation}
  \boxed{
  [\hat{r}_\alpha,\Vnl]|\psi\rangle
  = \sum_{I,l,m} h^l_I\!\left[
    |r_\alpha\beta^{lm}_I\rangle\langle\beta^{lm}_I|\psi\rangle
  - |\beta^{lm}_I\rangle\langle r_\alpha\beta^{lm}_I|\psi\rangle
    \right].
  }
  \label{eq:comm-result}
\end{equation}
This is a rank-2 update (in the language of matrix algebra): the
commutator of a separable operator with $\hat{r}_\alpha$ is
again separable, but now involves two types of projectors — one
weighted by $r_\alpha$ in the ket, one in the bra.

Equation~(\ref{eq:comm-result}) has a transparent structure:
\begin{itemize}
  \item The \emph{first term}
    $|r_\alpha\beta^{lm}_I\rangle\langle\beta^{lm}_I|\psi\rangle$
    keeps the unmodified projector in the bra and weights the ket
    by $r_\alpha$.
    Physically: the electron is projected onto the original KB
    projector at $\mathbf{r}'$ and then the result is multiplied
    by $r_\alpha$ at the field point $\mathbf{r}$.
  \item The \emph{second term}
    $|\beta^{lm}_I\rangle\langle r_\alpha\beta^{lm}_I|\psi\rangle$
    keeps the unmodified projector ket but weights the bra by
    $r_\alpha$.
    Physically: the electron is projected onto the
    $r_\alpha$-weighted projector (measuring overlap with the
    position-weighted projector) and the result propagates with
    the original projector shape.
\end{itemize}
The two terms have opposite signs and together ensure that the
commutator operator is anti-Hermitian, consistent with
$[\hat{r}_\alpha, \Vnl]^\dag = -[\hat{r}_\alpha, \Vnl]$ for
Hermitian $\hat{r}_\alpha$ and $\Vnl$.

\paragraph{Matrix element between Kohn-Sham orbitals.}
Taking the matrix element of eq.~(\ref{eq:comm-result}) between
$\langle\psi_a|$ and $|\psi_b\rangle$:
\begin{align}
  \langle\psi_a|[\hat{r}_\alpha,\Vnl]|\psi_b\rangle
  &= \sum_{I,l,m} h^l_I
     \Big[
       \underbrace{\langle\psi_a|r_\alpha\beta^{lm}_I\rangle}_{
         \displaystyle A^{lm}_{I,a}}
       \underbrace{\langle\beta^{lm}_I|\psi_b\rangle}_{
         \displaystyle B^{lm}_{I,b}}
     - \underbrace{\langle\psi_a|\beta^{lm}_I\rangle}_{
         \displaystyle \overline{B^{lm}_{I,a}}}
       \underbrace{\langle r_\alpha\beta^{lm}_I|\psi_b\rangle}_{
         \displaystyle \overline{A^{lm}_{I,b}}}
     \Big],
  \label{eq:comm-matrix-element}
\end{align}
where we introduced the notation
\begin{align}
  B^{lm}_{I,b}
  &= \langle\beta^{lm}_I|\psi_b\rangle
   = \int \beta^{lm*}_I(\mathbf{r})\,\psi_b(\mathbf{r})
     \,\mathrm{d}^3r,
  \label{eq:B-overlap}
  \\
  A^{lm}_{I,b}
  &= \langle r_\alpha\beta^{lm}_I|\psi_b\rangle
   = \int r_\alpha\,\beta^{lm*}_I(\mathbf{r})\,\psi_b(\mathbf{r})
     \,\mathrm{d}^3r.
  \label{eq:A-overlap}
\end{align}
Here $A^{lm}_{I,b}$ is the overlap of the $\alpha$-position-weighted
projector $r_\alpha\beta^{lm}_I$ with the orbital $\psi_b$, and
$B^{lm}_{I,b}$ is the standard KB projector overlap.
Both are computed as spatial integrals over the real-space grid.

Equation~(\ref{eq:comm-matrix-element}) therefore requires two types
of integrals per projector per orbital: the standard KB overlap
$B^{lm}_{I,b}$ (already available from the usual nonlocal PP force
evaluation) and the new position-weighted overlap $A^{lm}_{I,b}$.

%--------------------------------------------------------------------
\subsection*{C.\; Generalization to Multiple Projectors per Channel}
\label{app:multi-projector}
%--------------------------------------------------------------------

In the general KB form with multiple projectors per angular momentum
channel (eq.~11 of Kleinman and Bylander\cite{kleinman-bylander-1982}),
\begin{equation}
  \Vnl = \sum_I\sum_{l,m}\sum_{i,j}
  |\beta^{lm}_{I,i}\rangle\, h^l_{I,ij}\,
  \langle\beta^{lm}_{I,j}|,
  \label{eq:KB-general}
\end{equation}
where $i,j$ are radial projector indices and $h^l_{I,ij}$ is the
(generally non-diagonal) strength matrix.
The commutator result of eq.~(\ref{eq:comm-result}) extends directly
to
\begin{equation}
  [\hat{r}_\alpha,\Vnl]|\psi\rangle
  = \sum_{I,l,m}\sum_{i,j} h^l_{I,ij}\!\left[
    |r_\alpha\beta^{lm}_{I,i}\rangle
    \langle\beta^{lm}_{I,j}|\psi\rangle
  - |\beta^{lm}_{I,i}\rangle
    \langle r_\alpha\beta^{lm}_{I,j}|\psi\rangle
    \right].
  \label{eq:comm-result-general}
\end{equation}
The algebraic structure is identical; only the matrix of strength
coefficients $h^l_{I,ij}$ becomes non-diagonal in the projector
indices.
For the norm-conserving pseudopotentials used in \texttt{RMG}, the
single-projector case (eq.~\ref{eq:comm-result}) applies.

%--------------------------------------------------------------------
\subsection*{D.\; Nonlocal Contribution to the Current Density}
\label{app:Jnl-expression}
%--------------------------------------------------------------------

The nonlocal current density contribution to
eq.~(\ref{eq:J-total}) is defined through
eq.~(\ref{eq:Jnl-continuity}).
We now verify that this definition is consistent with the velocity
operator result by showing that the volume integral of
$\mathbf{J}_\mathrm{nl}$ gives the expected macroscopic contribution.

Using the single-particle density matrix $P(\mathbf{r},\mathbf{r}')
= \sum_n f_n \psi_n(\mathbf{r})\psi_n^*(\mathbf{r}')$, the
nonlocal contribution to the macroscopic current along direction
$\alpha$ is
\begin{align}
  J_{\mathrm{nl},\alpha}(t)
  &= \int_\Omega\!
     \operatorname{Re}\!\left[
       \psi^*(\mathbf{r},t)\,
       \frac{1}{\imath}[\hat{r}_\alpha,\Vnl]\,
       \psi(\mathbf{r},t)
     \right]\mathrm{d}^3r
  \nonumber \\
  &= \frac{1}{\imath}\sum_{I,l,m} h^l_I\,
     \operatorname{Re}\!\Bigg[
       \int\!\psi^*(\mathbf{r})\,r_\alpha\,\beta^{lm}_I(\mathbf{r})
       \,\mathrm{d}^3r
       \cdot
       \int\!\beta^{lm*}_I(\mathbf{r}')\,\psi(\mathbf{r}')
       \,\mathrm{d}^3r'
  \nonumber \\
  &\qquad\qquad\qquad\quad
   - \int\!\psi^*(\mathbf{r})\,\beta^{lm}_I(\mathbf{r})
     \,\mathrm{d}^3r
     \cdot
     \int\!r_\alpha\,\beta^{lm*}_I(\mathbf{r}')\,\psi(\mathbf{r}')
     \,\mathrm{d}^3r'
     \Bigg].
  \label{eq:Jnl-full}
\end{align}
In the density matrix notation of Section~\ref{sec:current-dm-sec},
this is evaluated as the matrix element
$\langle\psi_a|\frac{1}{\imath}[\hat{r}_\alpha,\Vnl]|\psi_b\rangle$
from eq.~(\ref{eq:comm-matrix-element}), contracted with $P_{ab}(t)$.
No additional real-space integration is required at each time step
beyond what is already computed in building $\Jmat{\alpha}_{ab}$.

%--------------------------------------------------------------------
\subsection*{E.\; Connection to \texttt{AppNls\_0xyz()} in RMG}
\label{app:code-connection}
%--------------------------------------------------------------------

In \texttt{RMG}, the commutator $[\hat{r}_\alpha,\Vnl]|\psi_b\rangle$
is evaluated by the function \texttt{AppNls\_0xyz()}, called from
\texttt{CurrentNlpp()} for directions $\alpha \in \{x,y,z\}$.

The function computes exactly eq.~(\ref{eq:comm-result}) on the
real-space grid.
For each orbital $\psi_b$ and each atom $I$, it evaluates:
\begin{enumerate}
  \item The standard projector overlap
    $B^{lm}_{I,b} = \langle\beta^{lm}_I|\psi_b\rangle$
    (a sum over grid points in the support of $\beta^{lm}_I$);
  \item The position-weighted overlap
    $A^{lm}_{I,b} = \langle r_\alpha\beta^{lm}_I|\psi_b\rangle$
    (the same sum but with each grid point contribution weighted
    by its $r_\alpha$ coordinate relative to atom $I$);
  \item The output array as
    $[\hat{r}_\alpha,\Vnl]\psi_b(\mathbf{r})
    = \sum_{I,l,m} h^l_I\!\left[
      r_\alpha\beta^{lm}_I(\mathbf{r})\cdot B^{lm}_{I,b}
      - \beta^{lm}_I(\mathbf{r})\cdot A^{lm}_{I,b}
    \right]$,
    stored in the work array \texttt{nv[]}.
\end{enumerate}
The result is passed to \texttt{CurrentNlpp()}, which computes the
matrix element $\langle\psi_a|\texttt{nv}\rangle$ via a BLAS
\texttt{gemm} operation, multiplies by \texttt{I\_t} $= \imath$,
and adds the result to \texttt{Palpha\_matrix}.
As shown in Section~\ref{sec:imath-convention}, the combined factor
$\imath \times \frac{1}{\imath} = 1$ ensures that the nonlocal PP
correction enters the total current matrix
$\boldsymbol{\mathcal{J}}^\alpha$ with unit weight, consistent with
eq.~(\ref{eq:Jmat-full}).

A self-consistency check: at the ground state, when all orbitals
$\psi_n$ are real-valued eigenstates, the density matrix is
$P_{ab} = f_a\delta_{ab}$ and the imaginary part of
$\langle\psi_a|\frac{1}{\imath}[\hat{r}_\alpha,\Vnl]|\psi_a\rangle$
should vanish because $[\hat{r}_\alpha,\Vnl]$ is anti-Hermitian
(the diagonal elements of an anti-Hermitian operator are purely
imaginary for complex orbitals and zero for real orbitals).
Thus $J_\mathrm{nl}(t=0) = 0$, consistent with vanishing ground-state
current.
This provides a simple numerical test that can be run on the
ground-state density matrix before starting the RT-TDDFT propagation.
